\section{Problem formulation and preliminaries}\label{sec:problem}
\subsection{Frozen prediction under recurring physical contexts}
Let $o_{e,t}$ and $a_{e,t}$ denote the observation and action at step $t$ of episode $e\in\{1,\ldots,N\}$. A frozen encoder and predictor produce
\begin{equation}
z_{e,t}=E(o_{e,t})\in\R^d,
\qquad \hat z^{\mathrm{base}}_{e,t+1}=f(z_{e,t},a_{e,t}).
\end{equation}
Each episode has an unobserved physical context $c_e$ that is fixed within the core task. Its dynamics satisfy
\begin{equation}\label{eq:dynamics}
z_{e,t+1}=f(z_{e,t},a_{e,t})+\mu_{c_e}(z_{e,t},a_{e,t})+\epsilon_{e,t},
\qquad \E[\epsilon_{e,t}\mid z_{e,t},a_{e,t},c_e]=0.
\end{equation}
Here $\mu_c$ is the conditional mean discrepancy from the frozen predictor, and the noise may depend on state and action. Context labels and physical parameter values are unavailable to the robot. Episode boundaries are observed, and data and memories may persist across them. Let $D_{<e}$ denote the observations, cues, actions, and transitions available before episode $e$.

The robot receives a contextual cue $y_e$ before its first informative interaction. Write $\tau_e$ for the decision step at which the corresponding action is committed; its context-dependent transition is observed afterward. The cue includes the context-relevant information available in the initial scene. For pre-contact motion history $H^{\mathrm{pre}}_e$, the core setting assumes
\begin{equation}\label{eq:noinfo}
I(c_e;H^{\mathrm{pre}}_e\mid y_e,D_{<e})=0.
\end{equation}
Thus moving to the interaction pose provides no additional identification of the current physics. This is a restriction on the information available in the task, not a claim that all robotic motion is uninformative.

\subsection{The first-contact limitation}
\begin{definition}[Single-contact task]
A task is single-contact with respect to an outcome component $J_e^{\mathrm{first}}$ if an action $a_e^{\mathrm{first}}$, or a fixed action segment, is committed at $\tau_e$ before feedback from its context-dependent transition is available, and later actions cannot change $J_e^{\mathrm{first}}$. The outcome component must be specified in advance and carry task cost or a safety consequence. A strict single-contact task has its entire task outcome determined this way.
\end{definition}
For puck striking, $J_e^{\mathrm{first}}$ can be the landing or stopping error after a single committed strike. For lifting, it can be peak acceleration during the initial force application; the later placement need not be irreversible. The definition concerns the commitment and feedback times, rather than literally counting all geometric contacts. A diagnostic nudge before the strike changes the information available and is excluded from the core protocol.

\begin{observation}[First-contact initialization]\label{obs:first}
Consider an adaptive predictor whose episode state $S_{e,t}$ changes only upon receipt of eligible current-episode prediction-error feedback, and whose prediction has no other adaptation path. If no such feedback is available before $a_e^{\mathrm{first}}$ is committed, then $S_{e,\tau_e^-}=S_{e,0}$. Consequently, its prediction for every given $(z,a)$ at that time equals its episode-initialized prediction. If the state is reset to the frozen predictor at each episode boundary, then
\begin{equation}\label{eq:firstidentity}
\hat z^{\mathrm{adapt}}_{e,\tau_e+1}(z,a)=f(z,a).
\end{equation}
\end{observation}
This identity is about prediction before feedback. Equality of task outcomes additionally requires the same planning rule, inputs, and randomization. The observation does not cover persistent initialization, cue-driven retrieval, auxiliary pre-contact updates, or informative probing. Moreover, Eq.~\eqref{eq:noinfo} alone does not imply that no prediction errors occur: the no-feedback condition in Observation~\ref{obs:first} must be enforced separately. Context-independent errors can change parameters without identifying the context.

The bridge to contextual adaptation is therefore specific. Previously acquired knowledge can improve initialization; pre-contact cues or predictive deployment history can select which knowledge to use. In the absence of both useful prior structure and informative evidence, arbitrary novel dynamics cannot be identified before interaction.

\subsection{Contextual prediction and composition}
The desired pre-contact belief and predictive distribution are
\begin{align}
\pi_{e,\tau_e^-}(c)&=p(c_e=c\mid y_e,D_{<e}),\label{eq:initialposterior}\\
p(z'\mid z,a,y_e,D_{<e})&=\sum_c\pi_{e,\tau_e^-}(c)\,p(z'\mid z,a,c,D_{<e}).\label{eq:mixture}
\end{align}
The mixture includes a predictive hypothesis for conditions not yet represented in memory. After contact, new transitions update both context belief and, where warranted, the selected dynamics memories. These are distinct operations: \emph{context inference} changes belief over existing hypotheses; \emph{parameter learning} changes a hypothesis's predictive distribution; and \emph{memory creation} adds a hypothesis for unexplained dynamics.

For $F$ physical factors, write $c=(c^1,\ldots,c^F)$ and model
\begin{equation}\label{eq:composition}
\mu_c(z,a)=\sum_{j=1}^{F}\mu^{j}_{c^j}(z,a)+\iota_c(z,a).
\end{equation}
The interaction term $\iota_c$ captures departures from additive factor effects. The shared prior in Section~\ref{sec:sharedprior} encodes this decomposition across complete-context memories. Prediction at an unseen combination transfers its additive component while retaining uncertainty about its unobserved interaction. We supply factor roles and cue regions as architectural structure, but no factor-value labels. Independent variation, adequate action excitation, and connected coverage of factor combinations constrain factor identification; unrestricted unsupervised disentanglement is not assumed.
